\documentclass{article}
\usepackage[backend=biber,natbib=true,style=alphabetic,maxbibnames=50]{biblatex}
\addbibresource{/home/nqbh/reference/bib.bib}
\usepackage[utf8]{vietnam}
\usepackage{tocloft}
\renewcommand{\cftsecleader}{\cftdotfill{\cftdotsep}}
\usepackage[colorlinks=true,linkcolor=blue,urlcolor=red,citecolor=magenta]{hyperref}
\usepackage{amsmath,amssymb,amsthm,enumitem,fancyvrb,float,graphicx,mathtools,soul,tikz}
\usetikzlibrary{angles,calc,intersections,matrix,patterns,quotes,shadings}
\allowdisplaybreaks
\newtheorem{assumption}{Assumption}
\newtheorem{baitoan}{Bài toán}
\newtheorem{cauhoi}{Câu hỏi}
\newtheorem{conjecture}{Conjecture}
\newtheorem{corollary}{Corollary}
\newtheorem{dangtoan}{Dạng toán}
\newtheorem{definition}{Definition}
\newtheorem{dinhly}{Định lý}
\newtheorem{dinhnghia}{Định nghĩa}
\newtheorem{example}{Example}
\newtheorem{ghichu}{Ghi chú}
\newtheorem{hequa}{Hệ quả}
\newtheorem{hypothesis}{Hypothesis}
\newtheorem{lemma}{Lemma}
\newtheorem{luuy}{Lưu ý}
\newtheorem{nhanxet}{Nhận xét}
\newtheorem{notation}{Notation}
\newtheorem{note}{Note}
\newtheorem{principle}{Principle}
\newtheorem{problem}{Problem}
\newtheorem{proposition}{Proposition}
\newtheorem{question}{Question}
\newtheorem{remark}{Remark}
\newtheorem{theorem}{Theorem}
\newtheorem{vidu}{Ví dụ}
\usepackage[margin=2cm]{geometry}
\def\labelitemii{$\circ$}
\DeclareRobustCommand{\divby}{%
    \mathrel{\vbox{\baselineskip.65ex\lineskiplimit0pt\hbox{.}\hbox{.}\hbox{.}}}%
}
\setlist[itemize]{leftmargin=5mm}
\setlist[enumerate]{leftmargin=*}

\title{Fundamentals of Programming: Final Term Exam 2026\\Đề Thi Cuối Kỳ Cơ Sở Lập Trình 2026}
\author{Nguyễn Quản Bá Hồng}
\date{\today}

\begin{document}
\maketitle

\begin{center}
    \Large
    Tổng quan đề thi
\end{center}

\begin{table}[H]
    \small
    \centering
    \begin{tabular}{|c|p{12cm}|c|c|c|}
        \hline
        {\bf No.} & {\bf Problem} & {\bf Time} & {\bf Memory} & {\bf Score} \\
        \hline
        1 & Quotient {\it\&} remainder in integer division -- Thương {\it\&} số dư trong phép chia số nguyên & 1 s & 1 GB & 100 \\
        \hline
        2 & Lengths of 3 sides {\it\&} perimeter of a non-degenerate triangle -- Độ dài 3 cạnh {\it\&} chu vi tam giác không suy biến & 1 s & 1 GB & 100 \\
        \hline
        3 & Perimeter of non-degenerate quadrilateral -- Chu vi tứ giác không suy biến & 1 s & 1 GB & 100 \\
        \hline
        4 & Numbers of evens {\it\&} odd in an interval -- Số số chẵn {\it\&} số số lẻ trên 1 đoạn số nguyên & 1 s & 1 GB & 100 \\
        \hline
        5 & Remainder of sum {\it\&} product of a sequence of integers -- Số dư của tổng {\it\&} tích của 1 dãy số nguyên & 1 s & 1 GB & 100 \\
        \hline
        6 & Sign of product of a sequence -- Dấu của tích của 1 dãy số & 1 s & 1 GB & 100 \\
        \hline
        7 & Prime {\it\&} composite -- Số nguyên tố {\it\&} hợp số & 1 s & 1 GB & 100 \\
        \hline
        8 & Number of primes {\it\&} composites in an interval -- Số số nguyên tố {\it\&} số hợp số trên 1 đoạn số nguyên & 1 s & 1 GB & 100 \\
        \hline
        9 & 1st-order linear recurrence -- Dãy truy hồi tuyến tính cấp 1 & 1 s & 1 GB & 100 \\
        \hline
        10 & 2nd-order linear recurrence -- Dãy truy hồi tuyến tính cấp 2 & 1 s & 1 GB & 100 \\
        \hline
    \end{tabular}
\end{table}
{\bf Lưu ý.} Thí sinh không cần phải pass hết bộ test mới được tính điểm, mà chỉ cần pass nhiều bộ test nhất có thể thì vẫn được tính điểm trên các bộ test đã được Online Judge chấp nhận là đúng.

\begin{baitoan}[Quotient {\it\&} remainder in integer division -- Thương {\it\&} số dư trong phép chia số nguyên]
	Cho 2 số nguyên $a,b$. Tìm thương $q$ {\it\&} số dư $r$ trong phép chia $a$ cho $b\ge1$, i.e., $a = bq + r$ với $0\le r\le b - 1$.
	\item {\sf Input.} Mỗi test gồm 1 dòng duy nhất chứa 2 số nguyên $a,b$ ($-2\cdot10^9\le a\le2\cdot10^9$, $0\le b\le2\cdot10^9$).
	\item {\sf Output.} Với mỗi test, nếu $b\ne0$ thì in ra thương $q$ {\it\&} số dư $r$ cách nhau 1 khoảng trắng, nếu $b = 0$ thì in ra $-1$.
	\item {\sf Sample.}
	\begin{table}[H]
		\centering
		\begin{tabular}{|l|l|}
			\hline
			\verb|quotient_remainder.inp| & \verb|quotient_remainder.out| \\
			\hline
			1 0 & -1\\
			\hline
			1 2 & 0 1 \\
			\hline
			-2 1 & -2 0\\
			\hline
		\end{tabular}
	\end{table}
\end{baitoan}

\begin{baitoan}[Lengths of 3 sides {\it\&} perimeter of a non-degenerate triangle -- Độ dài 3 cạnh {\it\&} chu vi tam giác không suy biến]
    Nhập vào 3 số nguyên $a,b,c\in\mathbb{Z}$. Kiểm tra xem chúng có phải là độ dài 3 cạnh 1 tam giác không suy biến không, biết
    \begin{equation*}
    	a,b,c\mbox{ là độ dài 3 cạnh 1 tam giác không suy biến}\Leftrightarrow\left\{\begin{split}
    		0 < a &< b + c,\\
    		0 < b &< c + a,\\
    		0 < c &< a + b,
    	\end{split}\right.
    \end{equation*}
    nếu phải, tính chu vi $P\coloneqq a + b + c$ của tam giác đó.
    \item {\sf Input.} Mỗi test gồm 1 dòng duy nhất chứa 3 số nguyên $a,b,c$ ($-2\cdot10^9\le a,b,c\le2\cdot10^9$).
    \item {\sf Output.} Với mỗi test, in ra $P$ nếu $a,b,c$ là độ dài 3 cạnh 1 tam giác không suy biến, in ra $-1$ nếu ngược lại.
    \item {\sf Sample.}
    \begin{table}[H]
        \centering
        \begin{tabular}{|l|l|}
            \hline
            \verb|triangle.inp| & \verb|triangle.out| \\
            \hline
            1 1 0 & -1\\
            \hline
            -1 1 1 & -1\\
            \hline
            1 1 1 & 3\\
            \hline
        \end{tabular}
    \end{table}
\end{baitoan}

\begin{baitoan}[Perimeter of non-degenerate quadrilateral -- Chu vi tứ giác không suy biến]
	Biết 4 số thực dương $a,b,c,d\in(0,\infty)$ là độ dài của 4 cạnh của 1 tứ giác khi {\it\&} chỉ khi số lớn nhất trong 4 số đó nhỏ hơn tổng của 3 số còn lại, i.e.,
	\begin{equation*}
		a,b,c,d\in(0,\infty)\mbox{ là độ dài 4 cạnh 1 tứ giác}\Leftrightarrow2\max\{a,b,c,d\} < a + b + c + d.
	\end{equation*}
	Tính chu vi của 1 hình tứ giác khi biết độ dài 4 cạnh của tứ giác đó.
    \item {\sf Input.} Mỗi test chứa gồm 1 dòng chứa 4 số nguyên $a,b,c,d$ ($-2\cdot10^9\le a,b,c,d\le2\cdot10^9$).
    \item {\sf Output.} Với mỗi test, nếu $a,b,c,d$ không là độ dài 4 cạnh của 1 tứ giác, in ra $-1$. Ngược lại, in ra 1 số nguyên là chu vi của tứ giác đã cho.
    \item {\sf Sample.}
    \begin{table}[H]
        \centering
        \begin{tabular}{|l|l|}
            \hline
            \verb|rectangle.inp| & \verb|rectangle.out| \\
            \hline
            1 1 1 1 & 4 \\
            \hline
            1 1 1 2 & 5\\
            \hline
            1 1 1 3 & -1\\
            \hline
        \end{tabular}
    \end{table}
\end{baitoan}

\begin{baitoan}[Numbers of evens {\it\&} odd in an interval -- Số số chẵn {\it\&} số số lẻ trên 1 đoạn số nguyên]
	Cho 2 số nguyên $a,b$. Tính số số chẵn {\it\&} số số lẻ trên đoạn $[a,b]\cap\mathbb{Z} = [a,a + 1,a + 2,\ldots,b - 1, b]$.
	\item {\sf Input.} Mỗi test gồm 1 dòng duy nhất chứa 2 số nguyên $a,b$ ($-2\cdot10^9\le a,b\le2\cdot10^9$).
	\item {\sf Output.} Với mỗi test, in ra số số chẵn {\it\&} số số lẻ trên đoạn $[a,b]\cap\mathbb{Z}$.
	\item {\sf Sample.}
	\begin{table}[H]
		\centering
		\begin{tabular}{|l|l|}
			\hline
			\verb|number_even_odd.inp| & \verb|number_even_odd.out| \\
			\hline
			0 1 & 1 1 \\
			\hline
			-1 1 & 1 2\\
			\hline
			-1 2 & 2 2\\
			\hline
		\end{tabular}
	\end{table}
\end{baitoan}

\begin{baitoan}[Remainder of sum {\it\&} product of a sequence of integers -- Số dư của tổng {\it\&} tích của 1 dãy số nguyên]
	Cho 1 dãy số nguyên $\{a_i\}_{i=1}^n = a_1,a_2,\ldots,a_n$, {\it\&} 1 số nguyên $b$, tìm số dư của tổng $\sum_{i=1}^n a_i$ {\it\&} của tích $\prod_{i=1}^n a_i = a_1a_2\ldots a_n$ khi chia cho $b$.
	\item {\sf Input.} Mỗi test gồm 2 dòng. Dòng 1 chứa 2 số nguyên $n,b$ ($1\le n\le2\cdot10^5$, $0\le b\le2\cdot10^5$). Dòng 2 chứa $n$ số $a_1,a_2,\ldots,a_n$ ($-2\cdot10^5\le a_i\le2\cdot10^5$, $\forall i\in[n]$).
	\item {\sf Output.} Với mỗi test, nếu $b = 0$ thì in ra $-1$. Nếu $b\ne0$, thì in ra 2 số dư
	\begin{equation*}
		\left(\sum_{i=1}^n a_i\right)({\rm mod}\ b) = (a_1 + a_2 + \cdots + a_n)({\rm mod}\ b),\ \left(\prod_{i=1}^n a_i\right)({\rm mod}\ b) = (a_1a_2\cdots a_n)({\rm mod}\ b),
	\end{equation*}
	cách nhau 1 khoảng trắng.
	\item {\sf Sample.}
	\begin{table}[H]
		\centering
		\begin{tabular}{|l|l|}
			\hline
			\verb|remainder_sum_product.inp| & \verb|remainder_sum_product.out| \\
			\hline
			1 0 & -1 \\
			1 & \\
			\hline
			2 1 & 0 0\\
			1 2 & \\
			\hline
			3 2 & 1 1\\
			1 -3 5 & \\
			\hline
		\end{tabular}
	\end{table}
\end{baitoan}

\begin{baitoan}[Sign of product of a sequence -- Dấu của tích của 1 dãy số]
    Nhập vào $n$ số nguyên. Tìm dấu của tích $n$ số nguyên đó với hàm dấu được định nghĩa bởi
    \begin{equation*}
    	{\rm sign}(x)\coloneqq\left\{\begin{split}
    		&1&&\mbox{if } x > 0,\\
    		&0&&\mbox{if } x = 0,\\
    		-&1&&\mbox{if } x < 0.
    	\end{split}\right.
    \end{equation*}
    \item {\sf Input.} Mỗi test gồm 2 dòng. Dòng 1 chứa duy nhất 1 số nguyên dương $n$ ($1\le n\le2\cdot10^5$). Dòng 2 chứa $n$ số nguyên $a_1,a_2,\ldots,a_n$ ($-2\cdot10^9\le a_i\le2\cdot10^9$, $\forall i = 1,2,\ldots,n$).
    \item {\sf Output.} Với mỗi test, in ra giá trị ${\rm sign}\left(\prod_{i=1}^n a_i\right) = {\rm sign}(a_1a_2\cdots a_n)$.
    \item {\sf Sample.}
    \begin{table}[H]
        \centering
        \begin{tabular}{|l|l|}
            \hline
            \verb|sign_product_sequence.inp| & \verb|sign_product_sequence.out| \\
            \hline
            3 & 0 \\
            0 1 2 & \\
            \hline
            4 & -1 \\
            -1 -2 -3 4 & \\
            \hline
        \end{tabular}
    \end{table}
\end{baitoan}

\begin{baitoan}[Prime {\it\&} composite -- Số nguyên tố {\it\&} hợp số]
	1 số nguyên dương $n\ge2$ được gọi là {\rm số nguyên tố} nếu nó có đúng 2 ước nguyên dương là $1$ {\it\&} $n$, {\it\&} $n$ được gọi là {\rm hợp số} nếu $n$ có nhiều hơn $2$ ước nguyên dương {\it\&} số ước của $n$ là hữu hạn, e.g., $2,3,5,7,11,\ldots$ là các số nguyên tố, trong khi $4,6,8,\ldots$ là hợp số. Kiểm tra xem 1 số là số nguyên tố hay hợp số, hay không là cả 2.
	\item {\sf Input.} Mỗi test gồm 1 dòng chứa duy nhất 1 số nguyên $n$ ($-2\cdot10^9\le n\le2\cdot10^9$).
	\item {\sf Output.} In ra chuỗi {\tt prime} nếu $n$ là số nguyên tố, in ra chuỗi {\tt composite} nếu $n$ là hợp số, nếu $n$ không phải số nguyên tố lẫn hợp số, in ra chuỗi {\tt none}.
	\item {\sf Sample.}
	\begin{table}[H]
		\centering
		\begin{tabular}{|l|l|}
			\hline
			\verb|prime_composite.inp| & \verb|prime_composite.out| \\
			\hline
			-2 & none\\
			\hline
			1 & none\\
			\hline
			2 & prime \\
			\hline
			4 & composite\\
			\hline
		\end{tabular}
	\end{table}
\end{baitoan}

\begin{baitoan}[Number of primes {\it\&} composites in an interval -- Số số nguyên tố {\it\&} số hợp số trên 1 đoạn số nguyên]
	Cho 2 số nguyên $a,b$. Tính số số nguyên tố, số hợp số, {\it\&} số số không là số nguyên tố cũng không là hợp số trên đoạn $[a,b]\cap\mathbb{Z} = [a,a + 1,a + 2,\ldots,b - 1, b]$.
	\item {\sf Input.} Mỗi test gồm 1 dòng chứa duy nhất 2 số nguyên $a,b$ ($-2\cdot10^9\le a,b\le2\cdot10^9$, $a\le b$, $b - a\le10^6$).
	\item {\sf Output.} Với mỗi test, in ra 3 số tự nhiên $p,c,n$ lần lượt là số số nguyên tố, số hợp số, {\it\&} số số không là số nguyên tố cũng không là hợp số trên đoạn $[a,b]$ cách nhau 1 khoảng trắng.
	\item {\sf Sample.}
	\begin{table}[H]
		\centering
		\begin{tabular}{|l|l|}
			\hline
			\verb|number_prime_composite.inp| & \verb|number_prime_composite.out| \\
			\hline
			1 2 & 1 0 1\\
			\hline
			0 3 & 2 0 2\\
			\hline
			0 4 & 2 1 2\\
			\hline
			-2 5 & 3 1 4\\
			\hline
		\end{tabular}
	\end{table}
\end{baitoan}

\begin{baitoan}[1st-order linear recurrence -- Dãy truy hồi tuyến tính cấp 1]
	Cho dãy số $\{a_n\}_{n=0}^\infty\subset\mathbb{Z}$ được xác định bởi công thức:
	\begin{equation*}
		\left\{\begin{split}
			a_0 &= \alpha\in\mathbb{Z},\\
			a_{n+1} &= aa_{n} + b,\ \forall n\in\mathbb{N}.
		\end{split}\right.
	\end{equation*}
	Tính $a_n$.
	\item {\sf Input.} Mỗi test gồm 1 dòng chứa 4 số nguyên $\alpha,a,b,n$ ($-2\cdot10^9\le\alpha,a,b\le2\cdot10^9$, $-2\cdot10^9\le n\le10^7$).
	\item {\sf Output.} Với mỗi test, Nếu $n < 0$ thì in ra $-1$. Nếu $n\ge0$ thì in ra giá trị của $a_n$ theo modulo $10^9 + 7$, i.e., $a_n({\rm mod}\ 10^9 + 7)$.
	\item {\sf Sample.}
	\begin{table}[H]
		\centering
		\begin{tabular}{|l|l|}
			\hline
			\verb|first_order_linear_recurrence.inp| & \verb|first_order_linear_recurrence.out| \\
			\hline
			0 1 1 2 & 2\\
			\hline
			0 1 2 2 & 4\\
			\hline
		\end{tabular}
	\end{table}
\end{baitoan}

\begin{baitoan}[2nd-order linear recurrence -- Dãy truy hồi tuyến tính cấp 2]
	Cho dãy số $\{a_n\}_{n=0}^\infty\subset\mathbb{Z}$ được xác định bởi công thức:
	\begin{equation*}
		\left\{\begin{split}
			a_0 &= \alpha,\ a_1 = \beta,\\
			a_{n+2} &= aa_{n+1} + ba_n + c,\ \forall n\in\mathbb{N},
		\end{split}\right.
	\end{equation*}
	Tính $a_n$.
	\item {\sf Input.} Mỗi test gồm 1 dòng chứa 6 số nguyên $\alpha,\beta,a,b,c,n$ ($-2\cdot10^9\le\alpha,\beta,a,b,c,n\le2\cdot10^9$, $-2\cdot10^9\le n\le10^7$).
	\item {\sf Output.} Với mỗi test, Nếu $n < 0$ thì in ra $-1$. Nếu $n\ge0$ thì in ra giá trị của $a_n$ theo modulo $10^9 + 7$, i.e., $a_n({\rm mod}\ 10^9 + 7)$.
	\item {\sf Sample.}
	\begin{table}[H]
		\centering
		\begin{tabular}{|l|l|}
			\hline
			\verb|second_order_linear_recurrence.inp| & \verb|second_order_linear_recurrence.out| \\
			\hline
			0 0 1 1 1 -2 & -1\\
			\hline
			0 0 1 1 1 1 & 0\\
			\hline
			0 0 1 1 1 2 & 1\\
			\hline
		\end{tabular}
	\end{table}
\end{baitoan}

%------------------------------------------------------------------------------%

\printbibliography[heading=bibintoc]

\end{document}